% GENERATED by gen_tables.py -- do not edit by hand (tab:surfaces)
\begin{table}[!htbp]
\caption{The five two-dimensional scans. $N$ is the total number of variational solves, which grows as the product of the two axis lengths; this quadratic cost is the reason the holdout experiment of Section~\ref{sec:holdout} matters.}
\label{tab:surfaces}
\centering\footnotesize\setlength{\tabcolsep}{3pt}
\adjustbox{max width=\columnwidth}{%
\begin{tabular}{llrr}
\toprule
Molecule & Grid (axis $\times$ axis) & $N$ & $n_q$ \\
\midrule
\ce{BeH2} & bond: 11 pts 0.9--2.4, angle: 10 pts 90--180 & 110 & 6 \\
\ce{NH3} & bond: 7 pts 0.85--1.6, angle: 7 pts 96--120 & 49 & 8 \\
\ce{H2O2} & dihedral: 10 pts 0--180, angle: 5 pts 80--140 & 50 & 10 \\
\ce{NO2} & bond: 11 pts 0.9--2.4, angle: 10 pts 90--180 & 110 & 6 \\
\ce{N2O} & rnn: 10 pts 0.9--2.25, rno: 6 pts 0.95--1.7 & 60 & 6 \\
\bottomrule
\end{tabular}}
\end{table}
